% DAC 2027 research manuscript format, checked against the official CFP on
% 2026-10-29. This is an explicitly hypothetical-results design draft.
% No invented numbers in this file may be promoted to measured evidence.
\documentclass[conference,9pt,letterpaper]{IEEEtran}
\usepackage[T1]{fontenc}
\usepackage{newtxtext,newtxmath}
\usepackage{graphicx}
\usepackage{booktabs}
\usepackage{amsmath}
\usepackage{array}
\usepackage{xcolor}
\usepackage{cite}
\usepackage{url}
\usepackage[hidelinks]{hyperref}
\usepackage{eso-pic}
\hypersetup{pdftitle={MECC: Single-GPU MoE Inference with Compiler-Guided Three-Tier Expert Caching -- Hypothetical Design Draft},pdfauthor={},pdfsubject={All evaluation outcomes are hypothetical design targets, not measurements. Not for submission.},pdfkeywords={MoE, single GPU, three-tier cache, hypothetical results}}
\AddToShipoutPictureFG{\AtPageLowerLeft{\put(45,18){\makebox(520,12){\normalfont\fontsize{8}{9}\selectfont\color{brown}HYPOTHETICAL RESULTS DESIGN DRAFT --- NOT MEASURED --- NOT FOR SUBMISSION}}}}
\newcommand{\mecc}{\textsc{MECC}}
\newcommand{\BH}{B_{\mathrm H}}
\newcommand{\BC}{B_{\mathrm C}}
\newcommand{\KH}{K_{\mathrm H}}
\newcommand{\KC}{K_{\mathrm C}}
\newcommand{\scenario}{\textit{Hypothetical evaluation.}}
\newtheorem{proposition}{Proposition}
\title{MECC: Single-GPU MoE Inference with\protect\\Compiler-Guided Three-Tier Expert Caching}
% Reserve a real author-block footprint for a later camera-ready version.
\author{\IEEEauthorblockN{Anonymous Authors}\IEEEauthorblockA{\mbox{}\\\mbox{}}}
\begin{document}
\maketitle
\begin{abstract}
Mixture-of-experts models activate few experts but retain a large weight footprint. Single-GPU deployment becomes storage-bound when expert weights exceed the task's host-memory allowance. We present MECC, a compiler-guided three-tier cache that combines profile-informed placement, cross-tier retention, and dependency-aware expert delivery without changing selected experts or weight precision. The compiler plans the initial state; runtime promotion and demotion adapt it to observed demand. In an explicitly hypothetical evaluation, MECC reduces single-GPU decode latency by $3.0$--$3.3\times$ over unoptimized streaming. A target configuration trades eight-GPU, one-second token latency for single-GPU, eight-second token latency, illustrating a deployment-cost tradeoff rather than multi-GPU speed parity.
\end{abstract}
\begin{IEEEkeywords}
Mixture of experts, single-GPU inference, expert offloading, three-tier caching, hardware/software co-design.
\end{IEEEkeywords}

\section{Introduction}
Mixture-of-experts (MoE) models increase parameter capacity while evaluating only a small subset of experts for each token~\cite{shazeer,mixtral}. This sparsity creates a promising deployment opportunity: one GPU can execute the selected experts even when the complete expert pool does not fit in its high-bandwidth memory (HBM). The opportunity is valuable for shared servers and smaller deployments, where assigning multiple accelerators to one inference task may be impractical. However, sparse computation does not eliminate the need to retain, locate, and deliver every expert that the router might select.

The relevant resource limit is the memory \emph{available to the inference task}, not the server's installed capacity. A server with substantial DRAM can still allocate a small pinned-memory allowance to one service. If HBM and the permitted host working set cannot retain the full expert pool, some selected weights must be fetched from SSD. The single GPU then alternates between useful computation and expensive weight delivery. Increasing host memory is one way to change this operating point; improving how a fixed allowance is used is the problem studied here.

Our goal is \emph{practical single-GPU deployment at a controlled memory budget}, rather than making one accelerator as fast as an arbitrarily large accelerator group. A useful evaluation therefore needs two complementary comparisons. An unoptimized single-GPU implementation reveals the benefit of our techniques under the same hardware constraints. A full-model-resident multi-GPU implementation identifies the latency paid to reduce the accelerator footprint. For example, one GPU taking eight seconds per token may be a valid deployment choice relative to eight GPUs taking one second, provided the quality and resource requirements are explicit. This example has equal GPU-seconds per token; it does not imply lower energy consumption or an eightfold reduction in execution cost.

Three challenges stand between sparse computation and this deployment goal. \textbf{First, fixed capacity must cover unequal demand.} Equal cache slots per layer are simple but may leave high-pressure layers repeatedly loading from SSD. \textbf{Second, useful experts must survive HBM eviction.} Independently managed tiers can retain redundant copies or evict an object shortly before its next use. \textbf{Third, a favorable initial placement must remain useful as requests change.} Offline popularity helps initialization, while runtime history is needed when the active working set moves. These are distinct effects and must be separated experimentally.

We propose \mecc{}, a compiler-guided expert-cache architecture spanning HBM, CPU pinned memory, and an SSD backing store. Its central idea is to use the profile as an initial prior and observed requests as online corrections. The compiler emits expert rankings and per-layer capacities under separate HBM and CPU allowances. The runtime promotes requested experts, retains displaced experts in the host tier, and schedules copies with explicit readiness dependencies. All router-selected experts execute at the original weight precision. Figure~\ref{fig:overview} illustrates the division of responsibility.

Our focus is coordinated deployment and isolated evaluation. The paper asks whether a compiled initial state, bounded retention, and efficient delivery produce a useful single-GPU latency--resource operating point after their costs are included.

We develop three contributions: (1) a budgeted expert-placement interface linking independent profiles to executable residency plans; (2) a runtime coordinating promotion, demotion, and copy readiness while preserving expert semantics; and (3) an evaluation separating single-GPU optimization gains, multi-GPU tradeoffs, and the effects of placement, retention, and overlap. All numerical outcomes below are \emph{hypothetical targets}, not completed experiments.

\begin{figure*}[t]
\centering\includegraphics[width=\textwidth]{figures/fig01_overview.pdf}
\caption{\mecc{} for single-GPU MoE deployment. Sparse router selections define demand, an independent profile determines initial residency, and runtime promotion/demotion changes the working set across HBM and bounded host memory. SSD remains the backing store. Expert identities and weights are unchanged. The diagram is schematic.}
\label{fig:overview}
\end{figure*}

\section{Deployment Setting and Challenges}
\subsection{Single-GPU Execution with Bounded Host Memory}
Consider a model with $L$ MoE layers. Layer $\ell$ contains $N_\ell$ experts, and each token selects $k_\ell$ experts. Each expert is an immutable weight object $e=(\ell,j)$ with size $s_e$. The GPU also holds the dense trunk, activations, KV cache, and temporary buffers. We use $\BH$ for the HBM allowance remaining for expert residency and $\BC$ for the \emph{task-level} CPU allowance. These are different from total physical capacities. Host-side mirrors and staging allocations must also be charged to the corresponding task budget.

The target regime has a persistent expert working set exceeding the permitted fast and intermediate storage. A nominal comparison of expert bytes with CPU bytes is only a deployment screen: the final configuration must additionally demonstrate genuine cold SSD reads after accounting for HBM residency. A workload that fits completely in HBM does not need this mechanism; one whose weights fit in host memory is a useful two-tier control, rather than proof of a three-tier benefit.

The execution preserves the router and its top-$k$ expert set. An HBM hit executes immediately. A host hit requires host-to-device (H2D) delivery before expert computation. An SSD miss adds an SSD-to-host read before H2D. We distinguish \emph{demand source}, observed when the request arrives, from the expert's location after prefetching. Otherwise, a prefetched SSD read could appear to be an HBM hit and obscure the actual storage work.

\subsection{Challenge 1: Place Limited Capacity Where It Matters}
Uniform layer allocation treats all layers as equally valuable cache consumers. Yet the benefit of another slot depends on the layer's request distribution and the cost of the miss it avoids. A layer with a small stable working set may require few slots; another may repeatedly visit experts outside its current coverage. The relevant question is not whether more memory helps, but whether \emph{the same total memory} can be placed more effectively.

Figure~\ref{fig:challenges}(a) specifies a controlled characterization: compare uniform and planned allocations while holding total slots, runtime replacement, precision, and trace fixed. The hypothetical outcome moves capacity toward layers with greater cold-read pressure. Some low-pressure layers may incur more misses after giving up slots; a useful plan reduces the aggregate exposed cost. This experiment must use a held-out request stream, since allocating from the evaluation histogram would leak future knowledge.

\subsection{Challenge 2: Preserve Reuse Across HBM Eviction}
An expert displaced from HBM may be used again before the system has visited many distinct experts. If the host tier is a fixed pool, the displaced expert may be discarded despite useful near-term reuse. If the two tiers retain duplicate objects independently, their combined unique coverage can also be smaller than their summed capacity. Both effects can increase SSD demand at the same nominal slot count.

Figure~\ref{fig:challenges}(b) contrasts a fixed host pool, independent LRUs, and coordinated retention. This is a comparison of complete residency rules, including their duplication behavior. It cannot by itself attribute the entire difference to recency updates. The follow-up analysis therefore records unique coverage, duplicate bytes, and migration traffic, and compares with an exclusive global-LRU reference. Figure~\ref{fig:retention} gives a concrete promotion/demotion example.

\subsection{Challenge 3: Combine a Useful Prior with Adaptation}
Profiling and online retention act at different time scales. A prior can avoid early misses before online history is available. Once the route distribution changes, however, continuing to pin yesterday's hot experts can waste capacity. Conversely, an online policy can adapt without a profile but must first discover the working set. A design that combines them should improve the finite serving horizon after initialization costs are included.

The appropriate experiment is factorial: profile versus random initialization crossed with static versus adaptive residency, followed by a no-shift control. Figure~\ref{fig:challenges}(c) illustrates the desired window-level observation. Importantly, its lower curve is not a universal property of joint LRU. Uniform accesses, weak locality, and inaccurate profiles may leave little benefit. The evaluation must retain those configurations and charge the cost of obtaining the prior.

\begin{figure*}[t]
\centering\includegraphics[width=\textwidth]{ideal_figures/fig02_challenges.pdf}
\caption{Three deployment challenges expressed as falsifiable experiments. (a) Same total slots, different layer allocations. (b) Same initial profile, different cross-tier retention rules. (c) Initialization and adaptation across a working-set shift. \textbf{All plotted values are hypothetical, not measured.} Layer groups and request windows are illustrative; no real route trace is claimed.}
\label{fig:challenges}
\end{figure*}

\section{MECC Design}
\subsection{Compiler: An Executable Residency Plan}
\mecc{} accepts model metadata, separate memory allowances, tier-delivery costs, and an independent route profile. For expert $e$, the profile gives request probability $p_e$. For a homogeneous layer, ranked probabilities can be summarized by a discrete Zipf model,
\begin{equation}
 p_{\ell,r}=\frac{r^{-\alpha_\ell}}{H_{N_\ell}(\alpha_\ell)},\qquad
 H_n(\alpha)=\sum_{r=1}^{n}r^{-\alpha}.
\end{equation}
Empirical expert identities remain attached to these ranks. The exponent summarizes concentration; it does not replace the identity map or encode request order. With enough observations, direct empirical probabilities can instead be used in the placement objective.

Let $x_{e,t}\in\{0,1\}$ indicate the expert's \emph{primary} tier, with $t\in\{\mathrm H,\mathrm C,\mathrm S\}$. Given delivery cost $T_t(s_e)$, the compiler models
\begin{align}
 \min_x\quad &\sum_e p_e\sum_t x_{e,t}T_t(s_e),\label{eq:placement}\\
 \text{s.t.}\quad&\sum_t x_{e,t}=1,\quad
 \sum_e s_e x_{e,t}\le B_t\quad(t\in\{\mathrm H,\mathrm C\}).\nonumber
\end{align}
An affine approximation $T_t(s)=a_t+s/b_t$ captures fixed overhead and effective bandwidth. The costs must be calibrated on the target machine; a calibration from another GPU is not an end-to-end prediction. The objective concerns delivery under a fixed prior. It does not model future locality, overlap perfectly, or certify runtime latency.

\begin{proposition}
For equal-size objects, fixed request probabilities, independent tier capacities, and $T_{\mathrm H}\le T_{\mathrm C}\le T_{\mathrm S}$, an optimal static placement assigns the most probable objects to HBM, the next most probable objects to CPU, and the remainder to SSD, filling available slots while objects remain.
\end{proposition}
\noindent\emph{Proof sketch.} If $p_i>p_j$ but object $i$ is assigned to a slower tier than object $j$, swapping their tiers changes expected cost by $(p_i-p_j)(T_{\mathrm{fast}}-T_{\mathrm{slow}})\le0$. Repeating these exchanges yields ranked placement. Using an empty faster-tier slot cannot increase cost. This argument applies to the static equal-size model with independent capacities; it does not establish optimality for heterogeneous expert sizes or a sequential multilayer heuristic.

For layer-specific capacities, the implementation allocates slots using marginal coverage increments. Under the Zipf summary, the next slot adds
\begin{equation}
 \Delta_{\ell}(r)=\frac{r^{-\alpha_\ell}}{H_{N_\ell}(\alpha_\ell)}.
\end{equation}
The current planner allocates HBM slots first and then CPU slots among the remaining ranked experts. This is an implementation heuristic for the joint multilayer problem. A cost-weighted objective can evaluate a proposed allocation, but the preceding proposition must not be stretched into a global optimality claim for that heuristic. Feasible plans require capacity checks, including any mandatory per-layer minimum.

The compiled artifact contains per-layer capacities, ordered expert IDs, initial placement, and modeled hit fractions. Model/profile identifiers and budget assumptions distinguish an independent profile from the natural-order fallback. The loader initializes residency directly from this artifact.

\subsection{Runtime: Coordinated Promotion and Retention}
For each layer, the runtime maintains ordered HBM and CPU residency maps. An HBM hit updates recency. A CPU hit removes the requested expert from the host residency set, transfers it to HBM, and updates recency there. An SSD miss first materializes the immutable weight object in a staging buffer, then transfers it to HBM. If the HBM map is full, its least-recently used victim is demoted into the CPU set; CPU pressure evicts the least-recently used host resident. SSD retains the backing object throughout.

Figure~\ref{fig:retention} illustrates the logical policy. Starting from HBM $[e_0,e_1]$ and CPU $[e_2,e_3,e_4]$, a request for $e_2$ promotes it and demotes $e_0$. The next request for $e_0$ can then be served without SSD. The illustration assumes exclusive logical residency and sufficient staging capacity. A real implementation may preserve a CPU mirror to avoid a device-to-host (D2H) copy; that mirror consumes physical bytes even though it is not a logical CPU-cache entry.

\begin{figure}[t]
\centering\includegraphics[width=\columnwidth]{figures/fig03_retention.pdf}
\caption{A concrete retention example, with both lists ordered LRU to MRU. Logical HBM capacity is two experts and CPU capacity is three. Promotion makes room in CPU for the HBM victim. This is a schematic state transition, not an experimental result. Physical mirrors and staging require additional explicit accounting.}
\label{fig:retention}
\end{figure}

Two distinctions matter for interpreting this runtime. First, exclusive joint retention is a coordination rule, not a newly invented general-purpose replacement algorithm. In the ideal object model used for controlled replay, its SSD demand-miss sequence matches a global LRU with capacity $\KH+\KC$ and a corresponding initial stack. Second, no policy dominates a static split on every request trace. A fixed pool can preserve a repeatedly useful expert that an adaptive policy evicts. We therefore evaluate observable workload conditions rather than assert unconditional superiority.

\subsection{Delivery: Overlap with Explicit Readiness}
The useful output of caching is an expert ready for computation, not merely an issued copy. A dedicated copy stream transfers pinned host weights. Each transferred expert receives a readiness event. The computation stream waits on that event before invoking the expert kernel, and the source buffer remains live until the copy completes. Eviction is delayed until all kernels using the resident weights have completed. These dependencies are required even when copy issuance and expert dispatch are separated by a generator or iterator interface.

The intended schedule overlaps delivery of a known upcoming expert with independent work. It does not infer future routes from evaluation labels. Prefetch depth, staging bytes, and predictive options are controlled separately. The primary replacement comparison disables predictive prefetch for every policy; the overlap ablation then enables the same known-demand transfer machinery for matched baselines. This separation prevents a faster transport path from being mistaken for a better replacement rule.

With unchanged weights, selected experts, kernels, and aggregation order, residency is a placement transformation of the original forward computation. Identity-preserving copies support semantic equivalence, but a byte-preserving design does not prove that an asynchronous implementation is correct. The intended correctness check compares synchronized reference execution and asynchronous execution on identical inputs, covering weight hashes, route selections, logits, generated tokens, and perplexity on a frozen text corpus.

\section{Implementation and Measurement Contract}
The prototype separates the compiler from a policy-driven engine. Model adapters map layer/expert identifiers to contiguous weight objects, which the SSD store serializes without changing precision. Residency policies share the loader, execution path, and instrumentation to isolate replacement effects. Small CPU state-machine checks complement GPU execution tests.

An implementable budget contract counts allocations by unique storage identity. GPU bytes include dense weights, resident experts, KV/activations, and transfer buffers. Host bytes include cache residents, mirrors of GPU residents, pinned staging, pageable buffers, and other process-owned weight copies. The ideal resource target in this draft is 70~GB GPU use under an 80~GB physical cap and 40~GB host use under a 40~GB task cap. The host target consists of 28~GB resident experts, 6~GB mirrors, 4~GB staging, and 2~GB other allocations. These figures are authored targets, not observed peaks.

Measurement windows are separated into initialization, warmup, prefill, decode, and quality evaluation. Each window has its own timer and traffic counters. Decode TPOT is elapsed decode time divided by the number of measured token intervals; TTFT includes prefill and the first output token but excludes model loading unless a cold-start result explicitly includes it. SSD demand source is recorded before any servicing action, while physical SSD/H2D/D2H bytes are recorded independently. Miss fraction, read bytes, and elapsed time answer different questions and should not be substituted for each other.

\section{Evaluation}
\scenario{} All numerical outcomes, repeat variation, quality checks, and resource peaks in this section are authored targets, not measurements.

\subsection{Experimental Methodology}
\textbf{Models and workloads.} The scenario uses Mixtral-8x7B and Mixtral-8x22B~\cite{mixtral,mixtral22}, each with question-answering and code prompts. Its single-A100 configuration has 80~GB physical HBM, a 50~GB nominal expert allowance subject to total-peak checks, and a 40~GB total task host-memory cap. Profile and evaluation splits are independent. Main comparisons use batch one, 128 input tokens, 128 output tokens, and matched deterministic decoding. Longer inputs, shorter outputs, and mixed-task sequences define sensitivity studies. Dataset and model revisions are execution artifacts to be frozen, rather than assumed available measurements.

\textbf{Unoptimized single-GPU baseline.} The primary baseline executes the same selected experts at FP16 precision on one GPU, with a fixed legal host staging allowance and demand-driven weight streaming. It has no profile-guided initial placement, dynamic CPU expert retention, or transfer/compute overlap. It remains a functioning out-of-core inference path: common kernels, tokenizer, model graph, and cold-read semantics are shared with \mecc{}. This baseline answers what the complete optimization stack changes relative to a straightforward single-GPU deployment. It is intentionally weak and cannot alone establish superiority over recent systems.

\textbf{Matched cache baselines.} A conventional three-tier LRU supplies a stronger reference. Local RF and SSD-native adapters represent recency/frequency replacement motivated by FlashMoE~\cite{flashmoe} and three-tier residency motivated by SSD-LLaMA~\cite{ssdllama}. They are hypothetical matched adapters, not official reproductions or published performance. MoE-CORE~\cite{core} is a particularly close prior/history comparator. These local comparisons do not establish state-of-the-art superiority over the external systems.

\textbf{Multi-GPU reference.} A resident implementation places the complete FP16 model across two, four, or eight A100s, where capacity permits. It uses the same model and inputs but no SSD expert delivery. It is a hardware-rich latency reference, not a same-budget baseline. GPU count, parallelism strategy, interconnect, software revision, batch size, and memory peak must accompany it. The illustrative eight-GPU reference is not attributed to FlashMoE or any external publication.

\textbf{Repetitions and metrics.} Three illustrative repeats at $0.97$, $1.00$, and $1.03$ times each mean define \emph{hypothetical spread}, not measured confidence intervals. Metrics include absolute TPOT, TTFT, throughput, demand-miss share, SSD/H2D/D2H bytes, and physical memory peaks. GPU-time comparisons fix the workload; energy and monetary cost require separate measurements.

\begin{figure*}[t]
\centering\includegraphics[width=\textwidth]{ideal_figures/fig04_performance.pdf}
\caption{\textbf{Hypothetical performance targets, not measurements.} (a) One A100 for every policy: the unoptimized baseline isolates the full optimization stack, while matched adapters supply stronger cache comparisons. Error bars encode an authored $\pm3\%$ spread. (b) An eight-A100 resident implementation provides a separate latency reference. Its hardware budget differs. RF and SSD-native bars are hypothetical local adapters, not reported results from the cited systems.}
\label{fig:performance}
\end{figure*}

\subsection{Single-GPU Performance}
Figure~\ref{fig:performance}(a) presents the intended comparison across four model/workload combinations. Hypothetical \mecc{} TPOT is 8.0, 9.5, 16.0, and 19.0~s, compared with 24, 31, 49, and 62~s for unoptimized streaming. The corresponding speedups are $3.00\times$, $3.26\times$, $3.06\times$, and $3.26\times$. These absolute values make the deployment latency visible; a normalized speedup alone would hide whether the single-GPU configuration is usable for the intended application.

The stronger SSD-native adapter takes 10, 13, 21, and 27~s per token in the scenario, corresponding to a more modest $1.25$--$1.42\times$ improvement for \mecc{}. This is the relevant scale for judging an incremental caching contribution. The large difference from the streaming baseline should not be used as the only headline. The hypothetical three-tier-LRU results also help separate the value of basic caching from the additional compiled initial state and delivery coordination.

The scenario TTFT values are 60/68/96/112~s for \mecc{} and 130/155/230/280~s for unoptimized streaming. Reporting both TTFT and TPOT exposes the prefill/decode tradeoff. Profile-driven warmup can improve decode while delaying first response, so we account for initialization separately rather than interpret decode speedup as end-to-end request speedup.

\subsection{Single GPU versus Multiple GPUs}
Figure~\ref{fig:performance}(b) separates the hardware-rich reference from single-GPU baselines. The illustrative eight-GPU TPOT values are 1.0/1.2/1.8/2.1~s. Thus the single-GPU design is slower in absolute token latency while using fewer accelerators. Whether that tradeoff is acceptable depends on the application's latency requirement and accelerator availability; the evaluation exposes the choice rather than declaring the single-GPU point universally preferable.

Table~\ref{tab:gpu} summarizes one deployment point. At one GPU and eight seconds per token, \mecc{} uses eight GPU-seconds per token. At eight GPUs and one second per token, the resident reference also uses eight GPU-seconds per token. The example demonstrates an eightfold reduction in allocated GPU count with an eightfold latency increase and equal nominal GPU-time. It does \emph{not} demonstrate an eightfold improvement in GPU efficiency. CPU time, SSD activity, utilization, interconnect power, and rental billing can change the actual operating cost.

\begin{table}[t]
\caption{Hypothetical 8x7B deployment point; not measured.}
\label{tab:gpu}\centering
\begin{tabular}{lrr}\toprule
Metric & Single-GPU \mecc{} & Resident reference\\\midrule
Allocated A100 GPUs & 1 & 8\\
Decode time / token (s) & 8 & 1\\
GPU-seconds / token & 8 & 8\\\bottomrule
\end{tabular}
\end{table}

\subsection{Attributing the Improvements}
Figure~\ref{fig:ablation} specifies an incremental ablation on the first workload. Adding profile placement to streaming changes hypothetical TPOT from 24 to 18~s; adding coordinated retention changes it to 11~s; enabling dependency-correct overlap changes it to 8~s. These are incremental stack comparisons. Because components interact, their differences are not independent additive contributions; a complete factorial experiment should also remove each component from the full system.

SSD demand-miss targets are 46\%, 33\%, 20\%, and 20\%. Placement and retention reduce cold demand; overlap preserves the miss count while reducing exposed waiting. This ablation targets the first workload. Figure~\ref{fig:challenges}(b) uses a separate controlled trace with a 12\% target.

The controlled comparisons separate these effects further. Placement compares profile and random initialization under the same runtime, including initial transfers. Retention compares fixed CPU, independent LRUs, and exclusive retention on one frozen trace, accounting for duplicates. Within-phase shuffling preserves expert histograms while removing short-range order. These controls distinguish frequency coverage from temporal locality; the authored shift remains a synthetic characterization, not an observation of natural MoE drift.

\begin{figure*}[t]
\centering\includegraphics[width=\textwidth]{ideal_figures/fig05_ablation.pdf}
\caption{\textbf{Hypothetical incremental-stack outcomes}, not completed ablations. (a) Absolute TPOT after adding each component. (b) Demand-source shares. The overlap-only step keeps admission, residency, and the demand-miss share unchanged; its intended benefit is reduced exposed transfer time.}
\label{fig:ablation}
\end{figure*}

\subsection{Budgets, Startup Cost, and Quality}
Figure~\ref{fig:budget}(a) sweeps the host allowance while fixing the GPU. The hypothetical \mecc{} TPOT decreases from 16 to 7~s as host budget increases from 16 to 48~GB; conventional LRU decreases from 27 to 12~s. The sweep identifies the value of better residency at each task allowance. Installed server capacity, assigned task cap, and physical use including mirrors and buffers remain distinct quantities.

Figure~\ref{fig:budget}(b) adds GPU count for the 8x7B question-answering workload. The one-GPU resident reference is infeasible. Two, four, and eight GPUs provide hypothetical resident TPOT of 2.8, 1.65, and 1.0~s. Corresponding \mecc{} targets are 8.0, 3.9, 2.1, and 1.25~s for one through eight GPUs. Including the smallest feasible resident configuration prevents an eight-GPU-only comparison from obscuring a cheaper alternative. These points describe a latency--resource frontier, not ideal linear scaling.

Figure~\ref{fig:budget}(c) charges the prior explicitly. With hypothetical initialization costs of 30~s for \mecc{} and 12~s for the SSD-native adapter, followed by 8 and 10~s per generated token, respectively, their total costs are equal after nine tokens. The first integer output length with strictly lower \mecc{} cost is ten. This simple illustration excludes TTFT and any profile cost not included in initialization, so a real end-to-end break-even analysis must either include those terms or label its scope precisely.

\begin{figure*}[t]
\centering\includegraphics[width=\textwidth]{ideal_figures/fig06_budget_cost.pdf}
\caption{\textbf{Hypothetical sensitivity targets.} (a) TPOT versus total task host-memory cap. (b) Single- and multi-GPU operating points; the one-GPU resident reference does not fit. (c) An initialization-plus-decode illustration includes an 18~s extra initial cost and a 2~s/token saving, giving equality at nine tokens. Model loading, profile generation, and prefill require explicit inclusion in an actual request-cost result.}
\label{fig:budget}
\end{figure*}

The scenario assumes identical weight hashes and selected experts, zero logit difference against deterministic synchronized execution, and zero perplexity difference on a fixed corpus. These are hypothetical pass conditions, not observed tests. Nondeterministic execution requires prespecified tolerances and token-level comparisons. Weight identity and stream readiness are separate obligations: preserving FP16 bytes does not itself establish correct asynchronous execution.

\section{Related Work and Scope}
\textbf{Expert offloading and heterogeneous execution.} Fast MoE offloading~\cite{offloading} exploits sparse expert execution under limited GPU memory. SwapMoE~\cite{swapmoe} studies tunable-memory serving; a local fixed-CPU adapter is not automatically a reproduction of its full published system. Fiddler~\cite{fiddler} coordinates CPU/GPU execution and uses popularity information in placement. Consequently, offline popularity cannot be claimed as absent from prior work. MoE-Lightning~\cite{lightning} combines scheduling and a performance model for memory-constrained inference, providing a relevant planning and transfer-scheduling reference.

\textbf{SSD-based MoE serving.} SSD-LLaMA~\cite{ssdllama} explicitly manages an SSD/RAM/VRAM hierarchy with dynamic expert delivery. FlashMoE~\cite{flashmoe} combines recency and frequency for replacement to reduce SSD I/O. Neither the existence of three tiers nor adaptive expert retention is unique to \mecc{}. Their native implementations and documented execution modes matter for comparisons; adapting their mechanisms into a common FP16 harness answers a narrower question than reproducing their published systems.

\textbf{Coordinated residency and lossless serving.} MoE-CORE~\cite{core} combines nonuniform layer capacities, domain-informed initialization, history-based replacement, and prefetching. ZipMoE~\cite{zipmoe} combines lossless compression with cache-affinity scheduling. These are close conceptual comparators. \mecc{} focuses on an executable placement artifact, explicit physical-budget accounting, and an isolated single-GPU deployment evaluation; the local-adapter results do not establish superiority over these external systems.

The scope is bounded-memory single-GPU serving. The GPU-count comparison establishes neither energy nor monetary savings. Hypothetical outcomes require replacement with measurements before submission.

\section{Conclusion}
\mecc{} coordinates profile-informed placement, cross-tier retention, and dependency-aware delivery for single-GPU MoE inference beyond the task's host-memory allowance. It preserves selected experts and precision. The hypothetical evaluation illustrates both the improvement over unoptimized single-GPU execution and the latency paid to reduce GPU count relative to resident deployment. Actual measurements must establish these operating points before submission.

\section*{Content Generated by AI}
OpenAI Codex (GPT-6) assisted in drafting all sections, generating figure code, and constructing hypothetical outcomes. These outcomes are not measurements. Human authors must verify all content, implementation correspondence, references, and replacement results and approve the final manuscript.

% DAC 2027: at most six BODY pages, plus one REFERENCES-ONLY page.
% Build check verifies the actual boundary; do not shrink the IEEE margins.
\clearpage
\IEEEtriggeratref{7}
\bibliographystyle{IEEEtran}
\bibliography{references}
\end{document}
